\documentclass[11pt,letterpaper]{article}

% ---------- Typography: serif body, sans headings (institutional memo standard) ----------
\usepackage[margin=1.05in,top=0.95in,bottom=1in]{geometry}
\usepackage[T1]{fontenc}
\usepackage{mathpazo}            % Palatino body — the classic finance-document serif
\usepackage{helvet}              % Helvetica for headings and labels only
\usepackage{microtype}           % refined kerning and justification
\usepackage{xcolor}
\usepackage[colorlinks=true,urlcolor=navy,linkcolor=navy]{hyperref}
\urlstyle{sf}                    % URLs in the document's sans font, not monospace
\usepackage{enumitem}
\usepackage{fancyhdr}
\usepackage{lastpage}
\usepackage{array}

\definecolor{navy}{HTML}{1F3864}
\definecolor{midgray}{HTML}{595959}
\definecolor{hairline}{HTML}{C8C8C8}

\linespread{1.10}                % comfortable executive leading
\setlength{\parindent}{0pt}
\setlength{\parskip}{3.5pt}

% ---------- Footer: understated, informational ----------
\pagestyle{fancy}
\fancyhf{}
\renewcommand{\headrulewidth}{0pt}
\fancyfoot[L]{\sffamily\fontsize{7.5}{9}\selectfont\color{midgray}ENNACIRI \textbar\ EARLY-WARNING MODEL FOR U.S. BANK CAPITAL DISTRESS \textbar\ JULY 2026}
\fancyfoot[R]{\sffamily\fontsize{7.5}{9}\selectfont\color{midgray}\thepage\ OF \pageref{LastPage}}

% ---------- Section label: small uppercase bold sans, navy — corporate memo register ----------
\newcommand{\sectionhead}[1]{%
  \vspace{6pt}%
  {\sffamily\bfseries\fontsize{9.5}{12}\selectfont\color{navy}\MakeUppercase{#1}}\par
  \vspace{-9pt}\textcolor{hairline}{\rule{\textwidth}{0.5pt}}\par\vspace{0pt}}

% Sub-label within a section
\newcommand{\sublabel}[1]{{\sffamily\bfseries\fontsize{9}{11}\selectfont\color{midgray}#1}\par\vspace{2pt}}

% Answer body
\newcommand{\answer}[1]{#1\par\vspace{3pt}}

\begin{document}

% ---------- Title block ----------
\vspace*{-14pt}
{\sffamily\fontsize{8.5}{11}\selectfont\color{midgray}MSDS 696 DATA SCIENCE PRACTICUM II \quad\textbar\quad WEEK 2 STATUS REPORT}\par\vspace{10pt}
{\sffamily\bfseries\fontsize{21}{25}\selectfont\color{navy}Early-Warning Model for\\U.S. Bank Capital Distress}\par\vspace{9pt}
{\sffamily\fontsize{9}{12}\selectfont Oussama Ennaciri \quad\textbar\quad Week 2 \quad\textbar\quad July 12, 2026}\par\vspace{5pt}
\textcolor{navy}{\rule{\textwidth}{1.2pt}}

% ---------- Project title ----------
\sectionhead{Project title}
\answer{Early-Warning Model for U.S.\ Bank Capital Distress}

% ---------- Project summary ----------
\sectionhead{Project summary}
\answer{This project builds an early-warning model that uses public FDIC and FRED quarterly data to flag which U.S.\ banks are at risk of a Prompt Corrective Action capital-tier drop. The data consists of five linked FDIC tables (1.68M bank-quarter records of each bank's quarterly financial and regulatory filings) plus a FRED macro table (economic context like interest rates, unemployment, and GDP). The next step is understanding the data to do feature selection.}

% ---------- Milestones ----------
\sectionhead{Milestones}
\answer{\begin{enumerate}[topsep=1pt,itemsep=1pt,leftmargin=16pt,label=\textcolor{navy}{\bfseries\arabic*.}]
  \item Literature review of bank failure and existing early-warning models (\textit{in progress})
  \item Collect the data (FDIC API, 1984 to present) (\textbf{DONE})
  \item Understand the data through exploratory data analysis (\textit{Next})
  \item Prepare the data for machine learning
  \item Train and compare three candidate models
\end{enumerate}}

% ---------- Last week's To Do ----------
\sectionhead{Last week's ``To Do''}
\answer{From the Week 1 proposal:
\begin{enumerate}[topsep=1pt,itemsep=1pt,leftmargin=16pt,label=\textcolor{navy}{\bfseries\arabic*.}]
  \item Literature review of bank failure and existing early-warning models (\textbf{in progress})
  \item Collect the data (FDIC API, 1992 to present)
  \item Understand the data through exploratory data analysis
  \item Prepare the data for machine learning
  \item Train and compare three candidate models
\end{enumerate}}

% ---------- This week's progress ----------
\sectionhead{This week's progress}
\answer{Built and ran a data pipeline via the FDIC and FRED APIs that pulls all five FDIC datasets (financials, institutions, failures, history) and the FRED macro table into working DataFrames: roughly 1.68M bank-quarter records spanning 1984 Q1 to 2026 Q1, joined against 175 quarters of macro context. Ran the data-fit checklist and confirmed the data has solid provenance, is feasible to acquire and use, and fits the question: it can directly detect bank distress, though it's not yet confirmed whether it can flag it early enough. Limitations noted: large datasets, quarterly lag, a rare target, and self-reported filings.}

% ---------- Issues & discussion ----------
\sectionhead{Issues \& discussion}
\answer{\begin{itemize}[topsep=1pt,itemsep=1pt,leftmargin=14pt,label=\textcolor{navy}{\footnotesize\textbullet}]
  \item The financials table has 2,378 fields per bank-quarter, too many to use directly. Narrowing this down to the right features requires time and further literature review.
  \item Becoming undercapitalized is rare, so the label is severely imbalanced. The evaluation metric needs to account for that.
  \item Call-report data is self-reported by banks, so a bank concealing problems can look healthy until it isn't.
  \item The data is filed quarterly, so it can miss fast, deposit-run-style collapses like SVB's, which happen between filings.
\end{itemize}
More limitations will likely surface as the data understanding and literature review continue.}

% ---------- Next week's To Do ----------
\sectionhead{Next week's ``To Do''}
\answer{\begin{enumerate}[topsep=1pt,itemsep=1pt,leftmargin=16pt,label=\textcolor{navy}{\bfseries\arabic*.}]
  \item Continuing the literature review
  \item Understanding the data to do feature selection
\end{enumerate}}

% ---------- Resources ----------
\sectionhead{Resources}
\answer{Literature review sources gathered so far:
\begin{itemize}[topsep=1pt,itemsep=1pt,leftmargin=14pt,label=\textcolor{navy}{\footnotesize\textbullet}]
  \item Cole and White, ``D\'{e}j\`{a} Vu All Over Again: The Causes of U.S.\ Commercial Bank Failures This Time Around,'' \textit{Journal of Financial Services Research}, 2012, Vol.\ 42, No.\ 1, pp.\ 5--29.
  \item Collier, Forbush, Nuxoll, and O'Keefe, ``The SCOR System of Off-Site Monitoring: Its Objectives, Functioning, and Performance,'' \textit{FDIC Banking Review}, 2003, Vol.\ 15, No.\ 3.
  \item Petropoulos, Siakoulis, Stavroulakis, and Vlachogiannakis, ``Predicting Bank Insolvencies Using Machine Learning Techniques,'' \textit{International Journal of Forecasting}, 2020, Vol.\ 36, No.\ 3, pp.\ 1092--1113.
  \item Correia, Luck, and Verner, ``Failing Banks,'' Federal Reserve Bank of New York Staff Report No.\ 1117, September 2024, revised June 2025.
  \item Chu, Frazier, Martin, Agarwal, Kuo, Northwood, Oey, Reed, Rodrigue, and Starr, ``Dissecting Depositor Flight: An Analysis of the Spring 2023 Bank Failures,'' FDIC.
\end{itemize}}

\vspace{2pt}
{\small\itshape\color{midgray}\textbf{Note on LLM use:} This document was typeset in LaTeX using Claude Code, which handled the formatting, styling code, and wording polish. All content, research, and decisions are my own, and every suggested edit was reviewed and approved by me.}

\end{document}
